\textbf{Resuelve} el siguiente problema de programación lineal mediante el algoritmo de \textit{Branch \& Bound} propuesto por Dakin indicando las cotas de cada nodo, cuáles son podados, etc:

\begin{equation}
\begin{split}
\mbox{max. } z = 15x_1 + 25x_2 +12x_3+10x_4\\
s.a.:\\
3x_1 + 6x_2 +5x_3 + 5x_4 \leq  12\\
x_1,\,\,x_2,\,\,x_3,\,\,x_4 \in  \{0,1\}
\end{split}
\end{equation}


Se propone \textbf{completar} este árbol en el que se indican, para cada subproblema, los valores de las variables $x_1$, $x_2$, $x_3$ y $x_4$ correspondiente a la solución óptima de su correspondiente relajación lineal.

\begin{figure}[h!]
    \begin{center}
    \begin{tikzpicture}[->,>=stealth',level/.style={sibling distance = 5cm/#1,
      level distance = 1.5cm}] 
    \node [arn_n, label=right: {(1, 1, 3/5, 0)}, label=above right: \textbf{\tiny{}}] {\tiny{$P_{0}$}}
        child{ node [arn_n, label=right: {(1, 1, 0, 3/5)}, label=below left: \textbf{\tiny{}}, label=above left: \textbf{\tiny{}}, label=below left: \textbf{\tiny{}}] {\tiny{$P_{1}$} }
            edge from parent node[above left] {$x_3 \leq 0$}
        	}
        child{ node [arn_n, label=right: {(1, 2/3, 1, 0)}, label=below right: \textbf{\tiny{}}, label=above right: \textbf{\tiny{}}] {\tiny{$P_{2}$}} 
            edge from parent node[above right] {$x_3 \geq 1$}
    		}
    ; 
    \end{tikzpicture}
    \end{center}
\end{figure}